% Part: first-order-logic
% Chapter: sequent-calculus
% Section: rules-and-proofs

\documentclass[../../../include/open-logic-section]{subfiles}

\begin{document}

\olfileid[id]{mvl}{seq}{rul}

\olsection{Aturan dan \usetoken{P}{derivation}}

Dalam pembahasan berikut, misalkan $\Gamma,\Delta,\Pi,\Lambda$
merepresentasikan barisan berhingga !!{sentence}s.

\begin{defn}[Sekuen]
Suatu \emph{sekuen bersisi~$n$} merupakan ekspresi berbentuk
\[
\Gamma_1 \nSequent \dots \nSequent \Gamma_n,
\]
dengan setiap $\Gamma_i$ merupakan barisan berhingga (mungkin kosong) yang
beranggotakan !!{sentence}s dari bahasa~$\Lang L$.
\end{defn}

\begin{defn}[Sekuen Awal]
Suatu \emph{sekuen awal bersisi~$n$} merupakan sekuen bersisi~$n$ berbentuk
$!A\nSequent\dots\nSequent !A$ bagi sebarang !!{sentence}~$!A$ dalam bahasa.

Jika bahasa tersebut memuat konektif $0$-tempat~$\star$, yakni suatu konstanta
proposisional, kita juga mengambil sekuen $\dots\nSequent\star\nSequent\dots$
dengan $\star$ muncul pada posisi bagi nilai kebenaran yang bersesuaian dengan
$\tf{\star}\in V$, dan posisi lainnya kosong.
\end{defn}

Bagi setiap konektif dari suatu logika bernilai~$n$, yakni~$\Log{L}$, terdapat
satu aturan logis bagi setiap nilai kebenaran yang mungkin diambil oleh
konektif tersebut dalam~$\Log{L}$. !!^{derivation}s dalam kalkulus sekuen
bersisi~$n$ bagi~$\Log{L}$ merupakan pohon sekuen: sekuen paling atas
merupakan sekuen awal; dan jika suatu sekuen berada di bawah satu atau lebih
sekuen lain, sekuen itu harus mengikuti sekuen di atasnya secara benar melalui
suatu aturan inferensi bagi konektif-konektif dari~$\Log{L}$.

\begin{defn}[Teorema]
!!^a{sentence}~$!A$ merupakan suatu \emph{teorema} dari logika
bernilai~$n$, yakni~$\Log{L}$, jika terdapat !!a{derivation} dari sekuen
bersisi~$n$ yang memuat~$!A$ pada setiap posisi yang bersesuaian dengan nilai
kebenaran yang ditetapkan dari~$\Log{L}$. Kita menulis
$\Proves[\Log{L}]!A$ jika $!A$ merupakan teorema dan
$\Proves/[\Log{L}]!A$ jika bukan.
\end{defn}

\begin{defn}[!!^{derivability}]
!!^a{sentence}~$!A$ \emph{!!{derivable} dari} suatu himpunan
!!{sentence}s~$\Gamma$ dalam logika bernilai~$n$, yakni~$\Log{L}$,
$\Gamma\Proves[\Log{L}]!A$, jika dan hanya jika terdapat suatu himpunan
bagian berhingga~$\Gamma_0\subseteq\Gamma$ dan suatu barisan~$\Gamma_0'$ yang
memuat !!{sentence}s dalam~$\Gamma_0$ sedemikian sehingga sekuen berikut
mempunyai !!a{derivation}:
\[
\Lambda_1\nSequent\dots\nSequent\Lambda_n.
\]
Di sini $\Lambda_i$ adalah~$!A$ jika posisi~$i$ bersesuaian dengan nilai
kebenaran yang ditetapkan, dan $\Gamma_0'$ dalam kasus lainnya. Jika $!A$
tidak !!{derivable} dari~$\Gamma$, kita menulis $\Gamma\Proves/ !A$.
\end{defn}

Sebagai contoh, logika \L ukasiewicz bernilai~$3$ mempunyai kalkulus sekuen
bersisi~$3$. Dalam sekuen bersisi~$3$,
$\Gamma\nSequent\Pi\nSequent\Delta$, $\Gamma$ bersesuaian dengan~$\False$,
$\Delta$ dengan~$\True$, dan $\Pi$ dengan~$\Undef$. Aksiomanya adalah
$!A\nSequent !A\nSequent !A$. Karena hanya $\True$ yang ditetapkan,
$\Gamma\Proves[\LogLuk[3]]!A$ jika dan hanya jika sekuen
$\Gamma\nSequent\Gamma\nSequent !A$ mempunyai !!a{derivation}. (Jika
$\Undef$ juga ditetapkan, kita memerlukan !!a{derivation} dari
$\Gamma\nSequent !A\nSequent !A$.)

\end{document}
